\documentclass[10pt,a4paper]{amsart}
\usepackage[margin=19mm]{geometry}
\usepackage{amsmath,amssymb,mathtools}
\usepackage[scr=rsfs,cal=euler]{mathalfa}
\usepackage{xfrac}
\usepackage[unicode,hidelinks,bookmarks=false]{hyperref}
\newcommand{\ie}{i.e.\ }
\newcommand{\cf}{cf.\ }
\newcommand{\resp}{respectively\ }
\newcommand{\Subsection}{\subsection}
\providecommand{\nobreakdash}{\nobreak\hbox{-}}
\setlength{\emergencystretch}{2em}
\multlinegap0pt
\usepackage{amssymb}
\newtheorem{mainthm}{Theorem}
\newtheorem{lem}{Lemma}
\newtheorem{prop}{Proposition}
\newtheorem{cor}{Corollary}
\newcommand{\A}{\mathfrak{A}}
\newcommand{\Tr}{\mathrm{tr}}
\newcommand{\X}{\mathcal{X}}
\newcommand{\bspi}{\boldsymbol{\pi}}
\newcommand{\dom}{\mathrm{dom}\,}
\newcommand{\frL}{\mathfrak{L}}
\newcommand{\g}{\gamma}
\renewcommand{\l}{\lambda}
\renewcommand{\phi}{\varphi}
\renewcommand{\rm}[1]{\mathrm{#1}}
\newcommand{\rmM}{\mathbf{M}}
\newcommand{\rmh}{\mathrm{h}}
\renewcommand{\t}[1]{\widetilde{#1}}
\newcommand{\tr}{\overline{\mathrm{tr}}}
\newcommand{\vol}{\mathrm{vol}}
\title{Entropy and determinants for unitary representations}
\author{Tim Austin}
\date{Source: arXiv:2412.13751v4 (CC BY 4.0)}
\begin{document}
\maketitle

\noindent\emph{Source and licence.} Entropy and determinants for unitary representations. Version arXiv:2412.13751v4 (CC BY 4.0), distributed by arXiv under the Creative Commons Attribution 4.0 International licence, http://creativecommons.org/licenses/by/4.0/, as recorded in the arXiv licence metadata for this exact version and shown on the abstract page https://arxiv.org/abs/2412.13751v4. Full licence notice: \texttt{licenses/arxiv-2412.13751-CC-BY-4.0.txt}.\par\smallskip

\noindent\textit{Editorial scope.} Counted unit: the article's mainthm mainthm:det-form (source0.tex lines 422--425). The article's complete original proof of it is delivered with it (source0.tex lines 2210--2247, one \texttt{proof} environment of 38 lines, which the article places away from the statement). No mathematics is written, repaired or re-derived: the statement and the proof are the article's own lines, in the article's own notation, and no retained line is substituted or re-flowed.  Retained uncounted as the source-local context the proof needs: lines 545--547; lines 553--565; lines 692--694; lines 763--765; lines 828--830; lines 854--862; lines 895--900; lines 1695--1705; lines 1750--1752; lines 1789--1791; lines 1971--1973; lines 1979--1983; lines 2001--2004; lines 2032--2043; lines 2145--2147; lines 2185--2187.  Nothing was omitted from any retained span, so the package carries no omission record.  No retained line contains a citation, so the wrapper carries no bibliography and no bibliography entry.  No retained line contains a figure environment, an image-inclusion command or drawing code, so no image is delivered and the package declares no asset dependency.  Editorial formatting only, in addition: an AMS \texttt{amsart} preamble that restates the article's own declarations and macros used by the retained lines, namely the environments \texttt{mainthm}, \texttt{lem}, \texttt{prop}, \texttt{cor} on separate counters, together with the standard packages those lines need.  The retained environments print in this document as Theorem 1, Lemma 1, Proposition 1, Corollary 1 in order of appearance, whereas the article prints the same items with the numbering of its own full document.  The complete native source of the article as distributed in the arXiv e-print for this exact version is bundled unchanged as \texttt{sources/170324/source0.tex}.
\begin{equation}\label{eq:y-Q-v}
[y_1,\dots,y_\ell]^\rm{T} := (I_{H_\pi}\otimes Q)[v_1,\dots,v_k]^\rm{T}.
\end{equation}

\begin{lem}\label{lem:new-type}
In the situation above, we have
\begin{equation}\label{eq:new-type}
\psi_{ii'}(b) = \sum_{j,j'=1}^k \phi_{jj'}((a_{ij})^\ast ba_{i'j'}) \qquad (b \in \A,\ 1\le i,i'\le k).
\end{equation}
As a result, with $[a_{ij}]$ held fixed, $\psi$ is continuous as a function of $\phi$ for the weak$^\ast$ topologies.

In particular, in the special case of~\eqref{eq:y-Q-v}, we have
	\begin{equation}\label{eq:psi-Q-phi}
\psi(b) := (Q^\rm{T})^\ast\phi(b)Q^\rm{T} \qquad (b \in \A).
	\end{equation}
\qed
\end{lem}

\begin{equation}\label{eq:pairing}
	\langle \phi,a\rangle := \frac{1}{k}\sum_{ij}\phi_{ij}(a_{ij}).
\end{equation}

\begin{prop}\label{prop:RadNik}
If $\phi$ is a $\l$-normal element of $\frL(\A,\rmM_k)_+$, then there is an operator $T$ affiliated to $\l^{\oplus k}(\A)'$ such that (i) $\xi_i \in \dom T$ for every $i$ and (ii) $\phi$ is associated to $\l^{\oplus k}$ by the tuple of vectors $T\xi_1$, \dots, $T\xi_k$.  Another affiliated operator $T_1$ also satisfies (i) and (ii) if and only if $|T_1| = |T|$.  In particular, the choice of $T$ is unique if we require it to be non-negative. \qed
\end{prop}

\begin{equation}\label{eq:det-phi-det-pairing}
\Delta_\tau \phi = (\Delta_{\t{\tau}\otimes \tr_k} T)^{2k} = \big(\Delta_{\tau\otimes \tr_k}(\langle \phi,\cdot\rangle)\big)^k.
\end{equation}

\begin{lem}\label{lem:Kap+2}
Let $k=1$, let $\phi_{\rm{ac}}$ be associated to $\l$ by $T\xi$ as in Proposition~\ref{prop:RadNik}, let $E_+$ be the orthogonal projection onto $(\ker T)^\perp$, and let $y := E_+\xi$.  Let $A$ be any dense $\ast$-subalgebra of $\A$. Then there is a sequence $(b_n)_{n\ge 1}$ of positive invertible elements of $A$ such that all of the following hold:
\begin{itemize}
\item[i.] $\tau(b_n(\cdot)b_n) \to \phi_{\rm{ac}}$;
\item[ii.] $\phi_{\rm{ac}}(b_n^{-1}(\cdot) b_n^{-1}) \to \Phi^\l_y$;
\item[iii.] $\phi_{\rm{sing}}(b_n^{-1}(\cdot) b_n^{-1}) \to 0$;
\item[iv.] $(\Delta_\tau b_n)^2 \to \Delta_\tau \phi$.
\end{itemize}
\end{lem}

\begin{prop}\label{prop:det-var-princ}
Let $A$ be any dense $\ast$-subalgebra of $\rmM_k(\A)$. Let $\tau$ be a tracial state on $\A$, and $\Delta$ its associated Fuglede--Kadison determinant.  Finally, let $\phi$ be an $\rmM_k$-valued completely positive map on $\A$.  Then
\begin{equation}\label{eq:det-var-princ}
(\Delta \phi)^{1/k} = \inf\big\{\langle \phi,a\rangle:\ a \in A\ \hbox{positive and invertible and}\ \Delta_{\tau\otimes \Tr_k}(a) \ge 1\big\}.
\end{equation}
\end{prop}

\begin{cor}\label{cor:sums}
Assume that $k=\ell=1$ and that $\phi$ and $\psi$ are disjoint.
\begin{enumerate}
\item[a.] For every neighbourhood $W$ of $\g$ there are neighbourhoods $U$ of $\phi$ and $V$ of $\psi$ such that
\[\X(\pi,W)\supset \X(\pi,U) + \X(\pi,V)\]
for any representation $\pi$.
\item[b.] For any neighbourhoods $U$ of $\phi$ and $V$ of $\psi$ there is a neighbourhood $W$ of $\g$ such that
\[\X(\pi,W) \subset \X(\pi,U) + \X(\pi,V)\]
for any representation $\pi$.
\end{enumerate}
\end{cor}

\begin{lem}\label{lem:strong-quot-for-pairing}
Let $(\pi_n)_{n\ge 1}$ be a sequence of separable representations, and let $\phi \in \frL(\A,\rmM_k)_+$.  Then $\phi$ is asymptotically associated to $(\pi_n)_{n\ge 1}$ if and only if $\langle \phi,\cdot\rangle$ is asymptotically associated to $(\pi_n^{(k)})_{n\ge 1}$. 
\end{lem}

\begin{lem}\label{lem:trace-conv-for-pairing}
Let $(\pi_n)_{n\ge 1}$ be an AP sequence for $\A$, and consider the AP sequence $(\pi^{(k)}_n)_{n\ge 1}$ for $\rmM_k(\A)$.  If $\tr_{d_n}\circ \pi_n\to \tau$, then $\tr_{kd_n}\circ \pi_n^{(k)}\to \tau \otimes \tr_k$.
\end{lem}

\begin{lem}\label{lem:upper-semicontinuity}
For any $\bspi$ and $k$, the function $\rmh_{\bspi}$ is upper semicontinuous on $\frL(\A,\rmM_k)_+$.
\end{lem}

\begin{lem}\label{lem:e-upper-bd}
Any $\phi \in \frL(\A,\rmM_k)_+$ satisfies
	\[\rmh_{\bspi}(\phi) \le \log \det \phi(1).\]
	In particular, if $\phi(1)$ is singular then $\rmh_{\bspi}(\phi) = -\infty$.
	\end{lem}

\begin{lem}\label{lem:AP-for-pairing}
Let $k$ be a positive integer, let $\bspi = (\pi_n)_{n\ge 1}$ be an AP sequence, and let $\bspi^{(k)} := (\pi_n^{(k)})_{n\ge 1}$.  Let $\phi\in\frL(\A,\rmM_k)_+$, and define $\langle \phi,\cdot\rangle$ as in equation~\eqref{eq:pairing}.  Then
\[\rmh_{\bspi^{(k)}}(\langle \phi,\cdot\rangle) = \frac{1}{k} \rmh_{\bspi}(\phi).\]
\end{lem}

\begin{prop}\label{prop:transform}
Fix $\phi \in \frL(\A,\rmM_k)_+$.
\begin{enumerate}
\item[a.] Let $Q \in \rmM_k$ be invertible, and define
\[\psi(b) := (Q^\rm{T})^\ast \phi(b) Q^\rm{T} \qquad (b \in \A).\]
Then
\[\rmh_{\bspi}(\psi) = 2\log |\det Q| + \rmh_{\bspi}(\phi).\]

\item[b.] Assume further that $\tr_{d_n}\circ \pi_n \to \tau$.  Let $a \in \rmM_k(\A)$ be invertible, and define $\psi \in \frL(\A,\rmM_k)_+$ in terms of $\phi$ and $a$ as in Lemma~\ref{lem:new-type}. Then $\phi$ is asymptotically associated to $\bspi$ if and only if $\psi$ is, and in that case
	\[\rmh_{\bspi}(\psi) = 2\log \Delta_{\tau\otimes \Tr_k}|a| + \rmh_{\bspi}(\phi).\]
\end{enumerate}
\end{prop}

\begin{cor}\label{cor:singular-part-controls}
Assume that $\tr_{d_n}\circ \pi_n \to \tau$.  Let $\phi \in \frL(\A,\rmM_k)_+$, and consider its Lebesgue decomposition $\phi_{\rm{ac}} + \phi_{\rm{sing}}$ relative to $\tau$.  Then $\phi$ is asymptotically associated to $\bspi$ if and only if $\phi_{\rm{sing}}$ is asymptotically associated to $\bspi$.
\end{cor}

\begin{cor}\label{cor:pre-B}
	Assume that $\tr_{d_n}\circ \pi_n \to \tau$.  Let $\tau$ be associated to $\l$ by $\xi$, let $a$ be a positive invertible element of $\A$, and let $\phi := \Phi^\l_{\l(a)\xi}$.  Then $\rmh_{\bspi}(\phi) = 2\log\Delta_\tau a$.
\end{cor}

\begin{mainthm}\label{mainthm:det-form}
Suppose that $d_n^{-1}\Tr_{d_n}\circ \pi_n \to \tau$ and that $\phi$ is asymptotically associated to $\bspi$. Then
\[\rmh_{\bspi}(\phi) = \log \Delta \phi_{\rm{ac}}.\]
\end{mainthm}

\begin{proof}[Proof of Theorem~\ref{mainthm:det-form}]
\emph{Step 1.}\quad First we prove the inequality ``$\ge$'' when $k=1$ and $\phi$ is $\l$-normal.  In this case parts (i) and (iv) of Lemma~\ref{lem:Kap+2} give a sequence of positive invertible elements $b_1$, $b_2$, \dots of $\A$ such that
\[\tau(b_i(\cdot)b_i) \to \phi \qquad \hbox{and} \qquad (\Delta b_i)^2 \to \Delta \phi \qquad \hbox{as}\ i\to\infty.\]
Then Corollary~\ref{cor:pre-B} gives $\rmh_{\bspi}(\phi_i) = 2\log \Delta b_i$ for each $i$. Letting $i\to \infty$, the upper semicontinuity from Lemma~\ref{lem:upper-semicontinuity} turns this into
\[\rmh_{\bspi}(\phi) \ge 2\lim_{i\to\infty}\log \Delta b_i = \log \Delta \phi.\]

\emph{Step 2.}\quad Next we prove the inequality ``$\ge$'' when $k=1$ but $\phi$ is otherwise arbitrary.  Let $O$ be any neighbourhood of $\phi$.  Since $\phi_{\rm{ac}}$ and $\phi_{\rm{sing}}$ are disjoint,  Corollary~\ref{cor:sums}(a) gives neighbourhoods $U$ of $\phi_{\rm{sing}}$ and $W$ of $\phi_{\rm{ac}}$ such that
\[	\X(\pi_n,O) \supset \X(\pi_n,U) + \X(\pi_n,W)\]
for every $n$.  Since $\phi$ is asymptotically associated to $\bspi$ by assumption, so is $\phi_{\rm{sing}}$ by Corollary~\ref{cor:singular-part-controls}.  Therefore $\X(\pi_n,U)$ is nonempty along an infinite subsequence of values of $n$, say $n_1 < n_2 < \dots$.  Let $\bspi'$ be the corresponding AP subsequence of $\bspi$.  For each $n_i$, the set $\X(\pi_{n_i},O)$ contains a translate of $\X(\pi_{n_i},W)$, and so
\begin{align*}
\limsup_{n\to\infty}\frac{1}{d_n}\log \frac{\vol_{2d_n}\X(\pi_n,O)}{v(d_n)} &\ge\limsup_{i\to\infty}\frac{1}{d_{n_i}}\log \frac{\vol_{2d_{n_i}}\X(\pi_{n_i},O)}{v(d_{n_i})}\\
&\ge \limsup_{i\to\infty}\frac{1}{d_{n_i}}\log \frac{\vol_{2d_{n_i}}\X(\pi_{n_i},W)}{v(d_{n_i})} \\
&\ge \rmh_{\bspi'}(\phi_{\rm{ac}}).
\end{align*}
This lower bound is at least $\log \Delta \phi_{\rm{ac}} = \log\Delta\phi$ by applying Step 1 along the AP sequence $\bspi'$.  Since $O$ is arbitrary, this proves that $\rmh_{\bspi}(\phi)\ge \log \Delta \phi$.

\vspace{7pt}

\emph{Step 3.}\quad We now prove the inequality ``$\le$'' in case $k=1$.  This proof is quickest via the variational principle from Proposition~\ref{prop:det-var-princ}.

Let $a \in \A$ be positive and invertible and satisfy $\Delta a \ge 1$.  Define a new positive functional by $\psi:= \phi(\sqrt{a}(\cdot)\sqrt{a})$.  Then we have
\begin{align*}
\log \phi(a) &= \log \psi(1) \\
&\ge \rmh_{\bspi}(\psi) &\qquad &\hbox{(Lemma~\ref{lem:e-upper-bd})}\\
&= 2\log \Delta \sqrt{a} + \rmh_{\bspi}(\phi) &\qquad &\hbox{(Proposition~\ref{prop:transform}(b))}\\
&\ge \rmh_{\bspi}(\phi) &\qquad &\hbox{(because $\Delta \sqrt{a} = \sqrt{\Delta a} \ge 1$)}.
\end{align*}
Taking the infimum over $a$, Proposition~\ref{prop:det-var-princ} turns this into $\log \Delta \phi \ge \rmh_{\bspi}(\phi)$.

\vspace{7pt}

\emph{Step 4.}\quad Finally, if $k > 1$, then we can apply the previous steps to the functional $\langle \phi,\cdot\rangle$ on $\rmM_k(\A)$.  First, Lemma~\ref{lem:trace-conv-for-pairing} gives that $\tr_{kd_n}\circ \pi_n^{(k)}\to \tau \otimes \tr_k$.  Secondly, Lemma~\ref{lem:strong-quot-for-pairing} gives that $\langle \phi,\cdot\rangle$ is asymptotically associated to $\bspi^{(k)}$.  Finally, we have
\begin{align*}
\rmh_{\bspi}(\phi) &= k\cdot \rmh_{\bspi^{(k)}}(\langle \phi,\cdot\rangle) &\qquad &\hbox{(Lemma~\ref{lem:AP-for-pairing})}\\
&= k\cdot \log \Delta_{\tau\otimes \tr_k}(\langle \phi,\cdot\rangle) &\qquad &\hbox{(case $k=1$ of Theorem~\ref{mainthm:det-form})}\\
&= \log \Delta \phi &\qquad &\hbox{(equation~\eqref{eq:det-phi-det-pairing})}.
\end{align*}
\end{proof}

\end{document}
